\documentclass[12pt]{article}

\setlength{\parindent}{0pt}
\setcounter{secnumdepth}{3}
\usepackage[margin=1in]{geometry}

% packages
\usepackage{amsmath, amsfonts, amsthm, amssymb}
\usepackage{graphicx, float}
\usepackage{mathtools}
\usepackage{titlesec}
\usepackage{interval}
\usepackage{hyperref}
\usepackage{siunitx}
\usepackage{titling}
\usepackage{vwcol}
\usepackage{setspace}
\usepackage{empheq}
\usepackage{cancel}
\usepackage{esdiff}
\usepackage{multicol}
\usepackage{mdframed}
\usepackage{esdiff}
\usepackage{tikzsymbols}
\usepackage{multicol}
\usepackage{tikz}
\usepackage{varwidth}
\usepackage{pgfplots}
\usepackage{fancyhdr}
\usepackage{enumitem}
\usepackage{tcolorbox}

\intervalconfig {
	soft open fences
}

% header
\pagestyle{fancy}
\fancyhf{}
\lhead{Ethan Fan}
\rhead{PCH}
\cfoot{\thepage}

% section format
\titleformat{\section}
  {\bfseries}
  {}
  {0em}
  {}
\titlespacing*{\section}{0pt}{0pt}{0.3em}

\titleformat{\subsection}[runin]
  {\bfseries}
  {}
  {0em}
  {}
\titlespacing*{\subsection}{0pt}{0pt}{0em}

\titleformat{\subsubsection}[runin]
  {\itshape}
  {(\roman{subsubsection})}
  {0em}
  {}

% commands
\newcommand{\alignedintertext}[1]{%
  \noalign{%
    \vskip\belowdisplayshortskip
    \vtop{\hsize=\linewidth#1\par
    \expandafter}%
    \expandafter\prevdepth\the\prevdepth
  }%
}

\newcommand*{\paren}[1]{\ensuremath\left(#1\right)}
\newcommand*{\limit}[2][x]{\ensuremath{\displaystyle\lim_{#1\to#2}}}
\newcommand*{\Limit}[3][x]{\ensuremath{\displaystyle\lim_{#1\to#2}\left[#3\right]}}
\newcommand*{\deriv}[1][x]{\ensuremath{\dfrac{\mathrm{d}}{\mathrm{d}#1}}}
\newcommand*{\Deriv}[2][x]{\ensuremath{\dfrac{\mathrm{d}}{\mathrm{d}#1}\left[#2\right]}}
\newcommand{\dydx}{\dfrac{dy}{dx}}
\newcommand*{\iinteg}[2][x]{\ensuremath{\displaystyle\int #2\;\mathrm{d}#1}}
\newcommand*{\dinteg}[4][x]{\ensuremath{\displaystyle\int_{#2}^{#3}#4\;\mathrm{d}#1}}
\newcommand*{\abs}[1]{\ensuremath{\left|#1\right|}}
\newcommand*{\eps}{\varepsilon}
\newcommand*{\real}{\ensuremath\mathbb{R}}
\newcommand*{\integer}{\ensuremath\mathbb{Z}}
\newcommand*{\posint}{\ensuremath\mathbb{Z^+}}
\newcommand*{\cbrt}[1]{\ensuremath{\sqrt[3]{#1}}}
\newcommand*{\lheqzero}{\ensuremath{\underset{\text{L'H}}{\overset{\left[\frac00\right]}{=}}}}
\newcommand*{\lheqinfty}{\ensuremath{\underset{\text{L'H}}{\overset{\left[\frac{\infty}{\infty}\right]}{=}}}}
\newcommand{\tor}{\text{ or }}
\newcommand{\floor}[1]{\lfloor #1 \rfloor}

\newcommand{\vem}{\vspace{0.5em}}
\newcommand{\example}[1]{\section{Example #1}}
\newcommand{\problem}[1]{\section{Problem #1}}
\newcommand{\subproblem}[1]{\subsection{(#1)} \,}

\DeclareMathOperator{\dne}{DNE}
\DeclareMathOperator{\sgn}{sgn}

\newcommand{\definition}[1]{\begin{tcolorbox}[colback=red!5!white,colframe=red!75!black,parbox=false] #1 \end{tcolorbox}}

\newcommand{\theorem}[2]{\begin{tcolorbox}[title={#1},colback=blue!5!white,colframe=blue!75!black,parbox=false] #2 \end{tcolorbox}}

\allowdisplaybreaks
% \postdisplaypenalty=100000

% document
\begin{document}

\vspace*{0cm}

% opening
\begin{center}
    {\Large \bfseries Problem Set 74} \\[0.5cm]
    {\large Derivatives of Inverse Functions} \\[0.35cm]
    {\normalsize May 8, 2026}
\end{center}

% problems

\problem{2}
Let $f(x) = x \arctan x$. Use the limit definition of derivative to find $f'(\sqrt{3})$.
\begin{align*}
    f'(\sqrt{3})
    &=\lim_{h \to 0}\frac{(\sqrt{3}+h)\arctan (\sqrt{3}+h)-\sqrt{3}\arctan \sqrt{3}}{h}\\
    &=\sqrt{3}\lim_{h \to 0} \frac{\arctan (\sqrt{3}+h)-\arctan \sqrt{3}}{h}+\lim_{h \to 0} \arctan     (\sqrt3+h)\\
    &=\sqrt{3}\lim_{h \to 0}\frac{\arctan \paren{\frac{\sqrt{3}+h-\sqrt{3}}{1+\sqrt{3}(\sqrt{3}+h)}}}{h}+\frac{\pi}{3}\\
    &=\sqrt{3}\lim_{h \to 0}\frac{\arctan \paren{\frac{h}{4+\sqrt{3}}}}{\paren{\frac{h}{4+\sqrt{3}}}}\cdot \frac{1}{4+\sqrt{3}h}+\frac{\pi}{3}\\
    &=\sqrt{3}\cdot 1 \cdot \frac{1}{4}+\frac{\pi}{3}\\
    &=\boxed{\frac{\sqrt{3}}{4}+\frac{\pi}{3}}
\end{align*}

\problem{3}
Suppose $f$ is a one-to-one differentiable function and its inverse function $f^{-1}$
is also differentiable. Use
implicit differentiation to show that
\[
(f^{-1})'(x)=\frac{1}{f'(f^{-1}(x))}
\]
provided that the denominator is not 0.
\begin{gather*}
    f(f^{-1}(x))=x\\
    \iff f'(f^{-1}(x))\cdot (f^{-1})'(x)=1\\
    \iff \boxed{(f^{-1})'(x)=\frac{1}{f'(f^{-1}(x))}}
\end{gather*}

\problem{6}
Let $f:\mathbb{R}\to \mathbb{R}$ be defined as $f(x)=x^5+e^{\frac{x}{2}}+e^{x^3}$ and $g(x)=f^{-1}(x)$. Find $g'(2)$.
\begin{gather*}
   2=x^5+e^{\frac{x}{2}}+e^{x^3} \implies  g(2)=0\\
   g'(2)=\frac{1}{f'(0)}=\frac{1}{\frac{1}{2}e^0+3\cdot0^2\cdot e^{0^3}}=\boxed{2}
\end{gather*}

\problem{7}
Find the derivatives of the following inverse trigonometric functions: \vem

\subproblem{e}
$y=\cot^{-1} x$
\begin{gather*}
    \cot y=x\\
    y'=-\frac{1}{\csc^2 x}\\
    y'=-\frac{1}{x^2+1}\\
    \boxed{\cot^{-1} x=-\frac{1}{x^2+1}}
\end{gather*}

\subproblem{f}
$y=\csc^{-1}x$
\begin{gather*}
    \csc y=x\\
    y'=-\frac{1}{\csc x\cot x}\\
    y'=-\frac{1}{|x|\sqrt{x^2-1}}\\
    \boxed{\csc^{-1}x=-\frac{1}{|x|\sqrt{x^2-1}}}
\end{gather*}

\problem{8}
Differentiate \vem

\subproblem{a}
$y = x^2 \arcsin x$
\begin{align*}
    \dydx&=\boxed{2x\arcsin x+\frac{x^2}{\sqrt{1-x^2}}}
\end{align*}

\subproblem{b}
$y = \dfrac{1 + \arctan x}{2 - 3 \arctan x}$
\begin{align*}
    \dydx 
    &= \frac{(2-3\arctan x)\frac{1}{x^2+1}-(1+\arctan x)(-\frac{3}{x^2+1})}{(2-3\arctan x)^2}\\
    &=\frac{\frac{1}{x^2+1}(2-3\arctan x+3+3\arctan x)}{(2-3\arctan x)^2}\\
    &=\boxed{\frac{5}{(x^2+1)(2-3\arctan x)^2}}
\end{align*}

\subproblem{c}
$f(x) = \sec^{-1} x \cdot \csc^{-1} x$
\begin{align*}
    \dydx
    &=\frac{\csc^{-1} x}{|x|\sqrt{x^2-1}}-\frac{\sec^{-1} x}{|x|\sqrt{x^2-1}}\\
    &=\boxed{\frac{\csc^{-1} x-\sec^{-1} x}{|x|\sqrt{x^2-1}}}
\end{align*}

\subproblem{d}
$f(x) = x(\tan^{-1} x)^2$
\begin{align*}
    \dydx
    &=(\arctan x)^2 + x \cdot 2\arctan x\cdot \frac{1}{x^2+1}\\
    &=\boxed{\arctan x(\arctan x+\frac{2x}{x^2+1})}
\end{align*}

\subproblem{e}
$f(x) = (\arcsin(x^3))^4$
\begin{align*}
    \dydx 
    &=4 (\arcsin (x^3))^3\cdot \frac{1}{\sqrt{1-x^6}}\cdot 3x^2\\
    &=\boxed{\frac{12x^2(\arcsin x^3)^3}{\sqrt{1-x^6}}}
\end{align*}

\subproblem{f}
$y = \arctan(e^{-x^2})$
\begin{align*}
    \dydx
    &=\frac{1}{e^{-2x^2}+1}\cdot e^{-x^2}\cdot (-2x)\\
    &=\boxed{-\frac{2xe^{-x^2}}{e^{-2x^2}+1}}
\end{align*}

\problem{9}
Let $y=\arctan \sqrt{\dfrac{a-x}{a+x}}, -a<x<a.$ Find $\dydx$. \\

Hint: Use substitution: $x=a\cos \theta.$
\begin{gather*}
    x=a\cos \theta \implies \theta =\arccos \frac{x}{a}\\
    y=\arctan\sqrt{\frac{a-x}{a+x}}=\arctan\sqrt{\frac{a(1-\cos \theta)}{a(1+\cos \theta)}}=\arctan (\tan \frac{\theta}{2})=\frac{\theta}{2}\\
    \dydx=\frac{1}{2}\cdot \frac{d\theta}{dx}=\frac{1}{2}(-\frac{1}{\sqrt{1-\frac{x^2}{a^2}}})(\frac{1}{a})=\boxed{-\frac{1}{2\sqrt{a^2-x^2}}}
\end{gather*}

\problem{11}
Find an equation of the line tangent to the graph of $y=\ln(x^2+4)-x\arctan \paren{\dfrac{x}{2}}$ at $x=2$.
\begin{gather*}
    f(2)=\ln (8)-2\arctan 1=3\ln 2-\frac{\pi}{2}\\
    f'(x)=\frac{2x}{x^2+4}-\arctan \paren{\frac{x}{2}}-x\cdot \frac{\frac{1}{2}}{1+\frac{x^2}{4}}\\
    f'(2)=\frac{1}{2}-\frac{\pi}{4}-\frac{1}{2}=-\frac{\pi}{4}\\
    \boxed{y-3\ln 2+\frac{\pi}{2}=-\frac{\pi}{4}(x-2)}
\end{gather*}

\problem{12}
Evaluate $\lim\limits_{h \to 0}\dfrac{(a+h)^2\arcsin (a+h)-a^2\arcsin a}{h}$.
\begin{align*}
    &\lim_{h \to 0}\dfrac{(a+h)^2\arcsin (a+h)-a^2\arcsin a}{h}\\
    =&\lim_{h \to 0}\dfrac{(a^2+2ah+h^2)\arcsin (a+h)-a^2\arcsin a}{h}\\
    =&\lim_{h \to 0}\dfrac{a^2(\arcsin (a+h)-\arcsin a)}{h}+\lim_{h \to 0}\dfrac{2ah\arcsin (a+h)}{h}+\lim_{h \to 0}\dfrac{h^2\arcsin (a+h)}{h}\\
    =&a^2\cdot \frac{1}{\sqrt{1-a^2}}+\lim_{h \to 0}2a\arcsin (a+h)+\lim_{h \to 0}h\arcsin (a+h)\\
    =&\boxed{\frac{a^2}{\sqrt{1-a^2}}+2a\arcsin a}
\end{align*}


\end{document}

